\begin{posterrow}{1,1}
  \postercolumn{
    \postersetup{accent=posterPlum}
    \begin{postersection}{When does staging data pay off?}
      Ensemble analyses often scan the same inputs with different parameters.
      Reading those inputs from shared storage on every pass consumes network
      bandwidth and can leave compute workers waiting \cite{example2025}.

      \begin{postercallout}{Research question}
        How much reuse does a workflow need before a node-local cache pays back
        the cost of copying the input?
      \end{postercallout}
      \postersubhead{Data paths}
      {\centering\resizebox{.74\linewidth}{!}{\datapaths}\par
      \postercaption{The cache fetches input from shared storage on the first pass.
        Later passes read the staged copy on the same node.}}
    \end{postersection}

    \posterspacer
    \postersetup{accent=posterBlue}
    \begin{postersection}{Benchmark design}
      Use the same input, worker count and analysis for both paths. Start each
      cached run with an empty cache and include the initial copy in wall time.

      \begin{posterrow}{1,1}
        \postercolumn{
          \postersubhead{What to measure}
          \begin{itemize}[nosep,leftmargin=*]
            \item End-to-end wall time, including staging.
            \item Bytes read from shared storage.
            \item Peak local storage per worker node.
          \end{itemize}
        }
        \postercolumn{
          \postersubhead{Test matrix}
          \begin{tabular*}{\linewidth}{@{}l@{\extracolsep{\fill}}l@{}}
            \hline
            Parameter & Setting \\
            \hline
            Input size & 64 GiB \\
            Worker nodes & 4 \\
            Passes & 1--4 \\
            Repeats per setting & 5 \\
            \hline
          \end{tabular*}
        }
      \end{posterrow}

      \posterspacer
      Alternate the run order and record shared-storage load to help distinguish
      caching effects from changes in contention.

      \textbf{Trade-off to report.} Four passes read 256 GiB from shared storage
      with direct reads versus 64 GiB with caching, assuming the staged copy
      stays resident. Report the peak local footprint alongside that reduction.
    \end{postersection}
  }
  \postercolumn{
    \postersetup{accent=posterTeal}
    \begin{postersection}{Runtime improves from the second pass}
      {\centering\resizebox{\linewidth}{!}{\runtimechart}\par
      \postercaption{Synthetic benchmark data. Points show illustrative medians.
        Error bars on the cached series show illustrative min--max ranges.}}
      The first cached pass costs 16\,s more than a direct read. By the fourth
      pass, staging saves 128\,s (31\%) in total wall time. The cached series
      grows more slowly because each additional pass reads from local storage.
    \end{postersection}

    \posterspacer
    \postersetup{accent=posterOchre}
    \begin{postersection}{A rule for deciding what to cache}
      With direct-read time $t_d$, cached-pass time $t_c$, and one-time staging
      cost $S$, a cache saves time after $n$ passes when
      \[ n(t_d-t_c)>S. \]
      In the synthetic example, $t_d=104\,\mathrm{s}$, $t_c=56\,\mathrm{s}$, and
      $S=64\,\mathrm{s}$, so two passes cover the staging cost.

      \begin{postercallout}{Where to test next}
        Vary input size and reuse together. Include runs where the working set
        exceeds local capacity to see how eviction changes the break-even point.
      \end{postercallout}

    \end{postersection}
  }
\end{posterrow}

\posterspacer
\begingroup
\postersetup{treatment=outline,accent=cleanborder,heading size=35bp}
\begin{posterrow}{1,1,1}
  \postercolumn{
    \begin{postersection}[equal height group=decisions]{One pass}
      \centering
      {\fontsize{52bp}{58bp}\selectfont\bfseries 16\,s slower\par}
      Staging costs 64\,s. One local read recovers only 48\,s of that cost.\par
    \end{postersection}
  }
  \postercolumn{
    \begin{postersection}[equal height group=decisions]{Two passes}
      \centering
      {\fontsize{52bp}{58bp}\selectfont\bfseries 32\,s saved\par}
      The second read pays for staging. Total time falls from 208\,s to 176\,s.\par
    \end{postersection}
  }
  \postercolumn{
    \begin{postersection}[equal height group=decisions]{Four passes with eviction}
      \centering
      {\fontsize{52bp}{58bp}\selectfont\bfseries 64\,s saved\par}
      A second staging step costs another 64\,s and halves the four-pass gain.\par
    \end{postersection}
  }
\end{posterrow}
\endgroup
